La Universidad del Valle de Guatemala se ha preocupado por el desarrollo de tecnología en esta área, particularmente empleando plataformas como el Crazyflie 2.1, con el objetivo de mejorar su eficiencia, seguridad y facilidad de uso. En 2025, Angela Figueroa \cite{figueroa2025desarrollo} y César Schwendener \cite{schwendener2025desarrollo}, estudiantes de la Universidad del Valle de Guatemala, estuvieron trabajando en la optimización de la infraestructura experimental y en herramientas de software para el control y evaluación del cuadricóptero Crazyflie 2.1. Sus resultados fueron satisfactorios, obteniendo una forma de controlar los drones utilizando \textit{Python} y reduciendo la latencia entre la computadora y el agente mediante mejoras en la comunicación. 

 

Debido a esta problemática, existen investigaciones relacionadas con el control de drones por medio de movimientos naturales, buscando una interacción más intuitiva. En 2016, Obaid et al. \cite{obaid2016would} realizaron un estudio centrado en el usuario para analizar cómo las personas podrían controlar drones mediante movimientos corporales. El objetivo fue definir un conjunto de acciones intuitivas para la navegación aérea. Más de 300 propuestas fueron recopiladas y analizadas, llegando a la conclusión de que existe un conjunto de comandos comunes en los cuales múltiples participantes coincidieron al realizar el mismo movimiento para representar acciones específicas como desplazamiento o rotación. Los resultados reflejaron que un  89\% de los gestos utilizados fueron dinámicos, a su vez el 45\% de las personas usan una mano para controlar un dron y el  42\% usaron ambas manos.

De manera similar,  se ha explorado la implementación física de estos sistemas de interacción corporal. En Suiza, Rognon et al. \cite{8304759} desarrollaron el exoesqueleto \textit{FlyJacket} para el control inmersivo de drones. Esta investigación tenía como objetivo hacer más accesible el control de drones mediante inclinaciones del torso, proporcionando una experiencia más natural para el usuario. Los resultados fueron satisfactorios en términos de inmersión y control; sin embargo, el principal inconveniente radicó en que el exoesqueleto dependía significativamente de la morfología del usuario y su ergonomía no resultó óptima para un uso prolongado, lo que limitó su aplicación práctica.

Otro estudio relevante fue la creación de un dispositivo portátil basado en una \textit{Raspberry Pi Zero W} y un sensor inercial de 6 ejes, el cual se coloca en la mano del piloto \cite{8662106}. Este sistema buscó clasificar diferentes poses de la mano y secuencias de movimientos que permitieran tanto un vuelo direccional como la ejecución de trayectorias predefinidas mediante determinadas acciones manuales. Los resultados fueron altamente satisfactorios, alcanzando una precisión del 98\% en el reconocimiento de comandos direccionales y del 92\% en la clasificación de trayectorias \cite{8662106}. Además, al evaluarse con usuarios con poca experiencia, el sistema demostró ser una alternativa más intuitiva para aprender a controlar un dron, así como una forma más interactiva y accesible de operación.

De manera complementaria, otros estudios han explorado el uso de sensores de captura de movimiento basados en visión para el control de drones, como el presentado en \textit{Gesture Control of a Drone Using a Motion Controller} \cite{suarez2016gesture}. En esta investigación, se desarrolló un sistema que permite controlar un dron mediante movimientos de la mano utilizando el sensor Leap Motion, el cual captura la posición y orientación de la mano en tiempo real. Para su implementación, se utilizó la plataforma \textit{Robot Operating System (ROS)} sobre un entorno \textit{Ubuntu}, junto con scripts desarrollados en \textit{Python} que procesaban la información del sensor y la convertían en comandos de vuelo para el dron. El sistema fue capaz de reconocer diferentes acciones como despegue, aterrizaje, navegación y acrobacias, mediante la interpretación de configuraciones específicas de la mano. Los resultados demostraron que el dron puede responder de manera estable y fluida a los movimientos del usuario, permitiendo un control intuitivo sin necesidad de entrenamiento avanzado. Asimismo, se evidenció que el sistema es flexible y escalable, ya que se pueden agregar nuevos comandos modificando los algoritmos de interpretación. Sin embargo, su funcionamiento depende de un área de interacción limitada y de condiciones controladas, lo que restringe su aplicabilidad en escenarios más amplios.

